% =====================================================================
%  2bat-ud3-14.tex  —  SESSIÓ 14: Repàs global i síntesi (estacions)
% =====================================================================
\horatitol{Sessió 14 — Repàs global i síntesi (estacions)}%
{Repàs actiu per quatre estacions, una per bloc, i mapa final que connecta límit, asímptota i continuïtat amb la lectura d'un model.}

\subsection{Idea clau}

Repassem \textbf{tota la unitat} de manera activa, rotant per \textbf{quatre
estacions}, una per cada bloc treballat. Al final construirem un \textbf{mapa de la
unitat} que connecta els tres conceptes —\textbf{límit}, \textbf{asímptota},
\textbf{continuïtat}— amb la \textbf{lectura global} d'un model.

\begin{idea}
\textbf{Les quatre estacions}
\begin{enumerate}[leftmargin=1.6em, itemsep=2pt]
  \item \textbf{Límits a l'infinit:} estimar el límit d'una funció des d'una taula o
        gràfica i interpretar-lo.
  \item \textbf{Asímptotes:} identificar les asímptotes d'una gràfica i dir-ne el
        significat social.
  \item \textbf{Continuïtat:} comprovar si una funció a trossos és contínua o fa un
        salt a una frontera.
  \item \textbf{Lectura global:} descriure el comportament complet d'un model social.
\end{enumerate}
\end{idea}

\subsection{Les estacions}

\begin{exemple}[Estació 1 — límits]
Per a $f(x)=500-\dfrac{1000}{x}$: a mesura que $x$ creix, $\dfrac{1000}{x}\to 0$, així
que $f(x)\to 500$. Per tant $\lim_{x\to\infty}f(x)=500$: el fenomen s'estabilitza en
$500$.
\end{exemple}

\begin{exemple}[Estació 3 — continuïtat]
$g(x)=\begin{cases}2x+1 & x\le 2\\ x+3 & x>2\end{cases}$ a $x=2$: esquerre
$2\cdot 2+1=5$; dret $2+3=5$. Coincideixen: \textbf{contínua}.
\end{exemple}

\subsection{Exercicis resolts}

\begin{exresolt}[model de l'Estació 2: asímptotes]
Per a $f(x)=\dfrac{20}{x}$, identifica les dues asímptotes i digues el significat de
cadascuna si $x$ és el nombre de beneficiaris d'un fons.
\solucio
\textbf{Horitzontal} $y=0$: quan $x$ creix, l'ajut per beneficiari tendeix a $0$ (es
dilueix amb molta gent). \textbf{Vertical} $x=0$: amb gairebé cap beneficiari, l'ajut
es dispara (situació impossible en el context).
\resposta{Horitzontal $y=0$ (l'ajut es dilueix); vertical $x=0$ (cap beneficiari, ajut
disparat).}
\end{exresolt}

\begin{exresolt}[model de l'Estació 4: lectura global]
Un model creix els primers mesos, s'acosta a una asímptota $y=300$ i al mes $15$ fa un
salt cap avall fins a $200$. Descriu-ne el comportament global.
\solucio
El servei \textbf{creix} al principi i s'\textbf{estabilitza} acostant-se a $300$
(asímptota: la seva capacitat). Al mes $15$ hi ha un \textbf{salt cap avall}
(discontinuïtat) fins a $200$: un canvi brusc (per exemple, una retallada) que redueix
de cop el servei. A partir d'aleshores es manté en $200$.
\resposta{Creix cap a $300$, salt a $200$ al mes $15$, i s'estabilitza en $200$.}
\end{exresolt}

\subsection{Exercicis per fer}

\noindent\textit{Una tasca per estació. Rota i resol-les totes.}

\begin{exercicis}
\begin{exlist}
  \item \textbf{(Estació 1)} Estima $\displaystyle\lim_{x\to\infty}\left(200-\frac{800}{x}\right)$
        fent una taula amb $x=10,100,1000$, i interpreta el resultat.
  \item \textbf{(Estació 2)} Per a $g(x)=150-\dfrac{300}{x}$, digues l'asímptota
        horitzontal i interpreta-la com l'ocupació d'un servei.
  \item \textbf{(Estació 3)} Comprova si
        $f(x)=\begin{cases}3x-2 & x\le 1\\ x & x>1\end{cases}$ és contínua a $x=1$.
  \item \textbf{(Estació 3)} Comprova si
        $h(x)=\begin{cases}x+4 & x\le 3\\ 2x & x>3\end{cases}$ és contínua a $x=3$ i,
        si no, digues la mida del salt.
  \item \textbf{(Estació 4)} Un model d'una llista d'espera creix i s'acosta a $y=600$
        sense cap salt. Descriu-ne el comportament global en dues frases.
  \item \textbf{(Mapa de la unitat)} Completa el mapa connectant cada concepte amb la
        seva idea:
        \begin{center}
        \renewcommand{\arraystretch}{1.4}
        \begin{tabular}{p{3.8cm} p{8.6cm}}
        \toprule
        \textbf{Concepte} & \textbf{Idea / lectura} \\
        \midrule
        Límit a l'infinit & \rule{7.8cm}{0.4pt} \\
        Asímptota horitzontal & \rule{7.8cm}{0.4pt} \\
        Asímptota vertical & \rule{7.8cm}{0.4pt} \\
        Continuïtat / salt & \rule{7.8cm}{0.4pt} \\
        \bottomrule
        \end{tabular}
        \end{center}
\end{exlist}
\end{exercicis}

% ---------------------------------------------------------------------
\fullsolucions{Sessió 14}
\begin{solucions}
\begin{sollist}
  \item $f(10)=120$, $f(100)=192$, $f(1000)=199,2$: s'acosta a $200$. Límit $200$; el
        fenomen s'estabilitza en $200$.
  \item Asímptota horitzontal $y=150$: el servei s'estabilitza en $150$ (per exemple,
        $150$ places), la seva capacitat.
  \item A $x=1$: esquerre $3\cdot 1-2=1$; dret $1$. Coincideixen: \textbf{contínua}.
  \item A $x=3$: esquerre $3+4=7$; dret $2\cdot 3=6$. Salt de mida $1$:
        \textbf{discontínua}.
  \item Resposta oberta. Idea: la llista creix al principi i s'estabilitza acostant-se a
        $600$ (el seu sostre); sense salts, l'evolució és suau i previsible.
  \item Resposta oberta. Idees: límit a l'infinit = valor cap al qual tendeix la funció
        quan $x$ creix; asímptota horitzontal = sostre de saturació ($y=L$); asímptota
        vertical = valor prohibit on la funció es dispara; continuïtat/salt = si el
        model canvia suaument o fa un canvi brusc.
\end{sollist}
\end{solucions}
